\documentclass{article}

\usepackage[dvipsnames]{xcolor}
\usepackage{listings}
\usepackage{amsmath}

\definecolor{JadeGreen}{RGB}{91,158,62}
\lstdefinestyle{catppuccin}{
	backgroundcolor=\color{White},
	commentstyle=\color{Gray},
	numberstyle=\footnotesize\ttfamily\color{Gray},
	stringstyle=\color{JadeGreen},
	keywordstyle=\color{BurntOrange},
	basicstyle=\ttfamily\footnotesize\color{Black},
	breakatwhitespace=false,
	breaklines=true,
	captionpos=b,
	keepspaces=true,
	numbers=left,
	numbersep=5pt,
	showspaces=false,
	showstringspaces=false,
	showtabs=false,
	tabsize=4,
}
\lstset{style=catppuccin}

\title{Software Saturday Lecture 2}
\author{Ryan Baker}
\date{\today}

\begin{document}

\maketitle
\tableofcontents
\pagebreak

\subsection*{Lecture Objectives}

\begin{itemize}
	\item Understand and use control structures (\texttt{if}, \texttt{switch}, loops)
	\item Define and call functions in C++
	\item Understand the concept of scope (local, global)
	\item Introduce arrays and basic array operations
	\item Use functions to produce maintainable code
\end{itemize}

\section{Control Structures}

	\subsection{\texttt{if} and \texttt{else} statements } \begin{itemize}
		\item The \texttt{if} statement is used to execute a block of code conditionally
		\item The \texttt{else} statement allows you to define an alternative
		\lstinputlisting[language=C++]{if.cpp}
		\item Use of \texttt{\{\}} to define the blocks (not needed  for 1-line blocks)
		\item Chaining statements together into \texttt{if () \{\} else if () \{\}}$\dots$
	\end{itemize}
	\subsection{\texttt{switch} statement} \begin{itemize}
		\item A \texttt{switch} statement evaluates and executes the corresponding \texttt{case}
		\item If no \texttt{case} matches, \texttt{default} is executed
		\lstinputlisting[language=C++]{switch.cpp}
		\item If the \texttt{break} statements aren't used, all blocks below will execute \begin{itemize}
			\item Sometimes, but rarely, this is desired
		\end{itemize}
	\end{itemize}
	
\subsection{Loops}

\subsubsection{\texttt{for} loops}

\begin{itemize}
	\item Used when you know how many times you need to repeat a block of code
	\item \texttt{// for (initialization; break condition; update) \{...\}}
	\lstinputlisting[language=C++]{for.cpp}
\end{itemize}

\subsubsection{\texttt{while} loops}

\begin{itemize}
	\item Used when you need a block to run until a condition is met
	\item \texttt{// while (condition) \{...\}}
	\lstinputlisting[language=C++]{while.cpp}
	\item \texttt{do \{\} while ()} loops \begin{itemize}
		\item Variation of \texttt{while} loops
		\item Runs the block first, performs condition check second
		\item As opposed to \texttt{while} which checks condition first and runs second
		\lstinputlisting[language=C++]{dowhile.cpp}
	\end{itemize}
\end{itemize}

\section{Functions}

\begin{itemize}
	\item Functions allow you to group statements into a reusable block
	\item This helps in reducing code repetition and organizing the program
	\item Defining functions: \texttt{type name(arguments) \{...\}} \begin{itemize}
		\item \texttt{type} is the return type of the function (\texttt{int}, \texttt{char}, \texttt{void}, ...)
		\item The function arguments are optional (can take 0 arguments)
		\item The function body \texttt{\{...\}} contains the statements that will run
	\end{itemize}
	\item Calling functions: \texttt{name(arguments)}
	\lstinputlisting[language=C++]{function.cpp}
	\item Function overloading \begin{itemize}
		\item C++ allows you to define multiple functions with the same name
		\item They must have different arguments
		\lstinputlisting[language=C++]{overload.cpp}
	\end{itemize}
\end{itemize}

\

\end{document}